\section{主线边界: 证据可追溯的科研智能体与数据要素服务}

\begin{designbox}{核心判断}
SciScope 的主线不是炫耀 AI Infra, 也不只是"又一个文献问答助手", 而是把公开科技文献
治理成\textbf{可计算、可调用的数据要素}, 并建成\textbf{可被其他科研 agent 调用的证据后端}。
大模型只是生成层; 可复现的数据资产、检索/证据索引、stance 判定与协议出口才是护城河。
\end{designbox}

\subsection{任务主线}

\begin{center}
\begin{tabularx}{\textwidth}{>{\bfseries}p{30mm}X X}
\toprule
主线模块 & 必须回答的问题 & 应形成的资产 \\
\midrule
数据要素治理 & 六源文献如何被规范化、去重、分区分层成可审计资产 & raw\_canonical / analysis / processed 分层 + 口径表 \\
数据分析报告 & 分布、演化、作者网络、覆盖偏差揭示了什么 & 图表、统计表、网络图、要素服务化收束 \\
证据立场 L3 & 一句话与文献是支持 / 反驳 / 中立, 证据落在哪个句子 & 金标准评测、校准/拒答、争议地图表与视图 \\
混合检索与问答 & 用户能否用自然语言得到带 `[n]` 证据的答案 & 混合检索、GraphRAG 证据问答 \\
MCP 协议出口 & 其他 agent 能否调用核查工具、读取争议资源 & MCP tools + resources、OpenCode 真调产物 \\
\end{tabularx}
\end{center}

\subsection{叙事比例（国赛答辩）}

\begin{plainbox}
\begin{itemize}
  \item \textbf{第一优先级（技术魂 A）}: 论断级 stance——不是"相关即支持", 而是
    证据句级支持/反驳/中立 + 校准拒答 + 限定条件 + 争议沉淀。
  \item \textbf{第二优先级（价值壳 C）}: 公开文献治理成数据要素, 以 API/MCP 对外服务,
    贴出题单位语境: 智能体作为文献要素服务的调用方。
  \item \textbf{第三优先级（展示皮 B）}: Go TUI 消费 SSE 的完整科研工作流, 本机全量可演示。
  \item \textbf{第四优先级}: 检索/趋势/推荐等赛题基线能力——必须能展示, 但不作为技术主叙事。
\end{itemize}
\end{plainbox}

\subsection{系统分层}

\begin{tikzpicture}[
  node distance=5mm,
  every node/.style={font=\small},
  layer/.style={draw=black, fill=white, rounded corners=2pt, minimum height=11mm, align=center, text width=0.88\textwidth},
  arrow/.style={-{Latex[length=2mm]}, thick, color=black}
]
  \node[layer, very thick] (mcp) {协议与出口层: API / SSE / MCP server（tools + resources）——被其他科研 agent 调用};
  \node[layer, below=of mcp] (evidence) {证据层（核心）: 混合检索, GraphRAG, 证据句定位, stance 判定, 校准/拒答, 争议地图};
  \node[layer, below=of evidence] (cap) {能力层: /api/search, /api/chat, /api/trends, /api/recommend, /api/graph, /api/agent/stream};
  \node[layer, below=of cap] (engine) {数据工程与模型层: src/harvest, src/analysis, src/models, src/infra, src/data\_contracts};
  \node[layer, below=of engine, densely dashed] (base) {可重建资产层: data/raw\_canonical, data/analysis, data/processed, models/, PostgreSQL/ pgvector};
  \draw[arrow] (base) -- (engine);
  \draw[arrow] (engine) -- (cap);
  \draw[arrow] (cap) -- (evidence);
  \draw[arrow] (evidence) -- (mcp);
\end{tikzpicture}

\clearpage